\section{Related Work}
\label{sec:Related_Works}

La letteratura recente nell'ambito della visione artificiale applicata alla medicina ha visto una convergenza tra architetture di classificazione profonda e modelli generativi.

\subsection{Classificazione e Residual Learning}

Il riconoscimento d'immagini moderno è dominato dalle Residual Networks introdotte da He et al. \cite{He_2016_CVPR}. La loro capacità di addestrare reti molto profonde tramite connessioni shortcut ha permesso di raggiungere performance d'eccellenza anche nel dominio medico. Nel contesto specifico delle radiografie toraciche, il lavoro di Kermany et al. \cite{kermany2018identifying} ha stabilito un benchmark fondamentale, dimostrando come modelli pre-addestrati possano assistere i clinici nella diagnosi della polmonite pediatrica.

\subsection{Evoluzione delle Generative Adversarial Networks}

Le GAN hanno rivoluzionato la generazione di dati sintetici. Le Deep Convolutional GAN (DCGAN) \cite{radford2015unsupervised} hanno introdotto vincoli architettonici basati su livelli convoluzionali che hanno reso possibile la generazione di immagini coerenti. Per permettere una generazione mirata alle classi specifiche, Mirza e Osindero \cite{mirza2014conditionalgenerativeadversarialnets} hanno proposto le Conditional GAN (cGAN), estendendo l'architettura per includere informazioni supplementari come input sia per il generatore che per il discriminatore.

\subsection{Addestramento Stabile e Wasserstein GAN}

Uno dei problemi storici delle GAN è l'instabilità dell'addestramento. La Wasserstein GAN (WGAN) \cite{arjovsky2017wasserstein} ha affrontato questo problema sostituendo la divergenza di Jensen-Shannon con la distanza di Earth-Mover, fornendo un gradiente più informativo. Per superare le limitazioni del weight clipping originariamente proposto, Gulrajani et al. \cite{gulrajani2017improved} hanno introdotto la Gradient Penalty (WGAN-GP), garantendo il rispetto del vincolo di Lipschitz in modo più fluido e robusto. L'unione di queste tecniche rappresenta l'attuale stato dell'arte per la generazione di immagini mediche ad alta fedeltà.

Le potenzialità di queste immagini sintetiche si estendono anche al superamento delle discrepanze tra dataset di diversa provenienza. [...continue DA...].
